\scanpage{36}
par récurrence sur $\operatorname{card}I$ la notion de
\unsure{donnée|données} \deleted{de} côniques sur \deleted{cette base}
$(X,(X_i)_{i\in I},(B_\alpha,\mathcal E_\alpha)_{\alpha\in\Lambda})$.
\deleted{On procède} Si $I=\varnothing$, ce sera la structure
vide (il y en a une et une seule). Si
$I\ne\varnothing$, on prend $I_0$ ens des éléments minimaux
de $I$, et une donnée cônique sera par
définition la donnée

\textcircled{a} d'un \deleted{tube}\inserted{vois. con. strict} $\overline T_{I_0}$
de \deleted{\unsure{la somme}} $X_{I_0}=\bigcup_{\beta\in I_0}X_\beta$,
transversal aux $(B_\alpha,\mathcal E_\alpha)$ et
\deleted{\unsure{donc} admissible}
côniquement transversal aux $X_i$ ($i\in I$).
--- ce qui permet donc d'introduire les
structures (19), et en particulier

\[
(22)\qquad
X^*,\quad
\underbrace{(X_i^*)_{i\in I^*=I\setminus I_0}}
_{\text{stratification de }I^*},
\quad
\underbrace{\big((\deleted{\unsure{\mathcal E}}B_\alpha,\mathcal E_\alpha)_{\alpha\in\Lambda},
(\dot X_\beta,\mathcal E_\beta)_{\beta\in I_0}\big)}
_{\substack{\text{famille transverse}\\\text{de sous-espaces}\\
\text{bordants ``cylindrisés''}\\
\text{de }X,\ \text{transversaux}\\\text{aux }X_i}}
\]

en plus de
\[
(23)\qquad
\big((X_\beta)_{\beta\in I_0},
(B_{\beta\alpha},\mathcal E_{\beta\alpha})_{\substack{\beta\in I_0\\\alpha\in\Lambda}}\big)
\]
On doit avoir de plus : $\unsure{\infty|a}$

\textcircled{b} Une structure cônique sur (22) (ce qui
a un sens, par récurrence\ldots)
\margin{35}
